\section{Arithmetic estimates for finite twist families}
\label{sec:transfer-criteria}

The pointwise converse will follow by placing the given curve at the
zero vertex of a large finite family of twists. We arrange analytic
nonvanishing at every other vertex, together with an upper bound on a
normalized valuation. The estimates in this section explain which
normalization is needed and what must be uniform as the family grows.
Their proofs occupy Sections~\ref{sec:binary-algebra},
\ref{sec:cy:interpolation} and~\ref{sec:heegner-interpolation}.

Fix $E/\Q$ of conductor $N$ and retain the local weights $w(h)$,
finite Sha lengths $s(h)$, and period-normalized values $\mathcal L(h)$
of Section~\ref{sec:preliminaries}. In particular $s(h)$ is defined at
a corank-zero base and at every twist already known to have analytic
rank zero or one. All valuations satisfy $v(2)=1$.

A binary family has the form
\[
 h_x=h_0\prod_{q\in\mathcal Q}(q^*)^{\lambda_q(x)},
 \qquad x\in\F_2^b,\qquad q^*=(-1/q)q,
\]
where the new primes are distinct and prime to $2Nh_0$, and each
$\lambda_q$ is linear. The old factors of $h_0$ remain at every vertex;
no new factor is present at zero. We will hold the required local
squareclasses fixed while $b$ grows. Both the set of primes and their
sizes may depend on $b$.

Here is the common transfer mechanism. For families satisfying the
hypotheses of the transfer theorems below, the arithmetic proofs attach
power series
\[
                 u_x(t)\in\Ztwo[[t]][1/2]
\]
to the vertices, using a common family of cochain models and compatible
trivializations of the determinant lines. The differences of these
coordinates therefore have a fixed meaning. Their constant terms detect
the central value in rank zero or a Heegner point in rank one. If the analytic nonvanishing sought
at the base fails, then $u_0(0)=0$. The nonzero vertices instead satisfy
$v(u_x(0))\le B_1$ with $B_1$ independent of $b$.
The denominator analysis proves that, for every fixed integer $m$ and
sufficiently large $b$, some $x\ne0$ satisfies
\[
                   v\bigl(u_x(0)-u_0(0)\bigr)\ge m.
\]
Taking $m>B_1$ contradicts these two assertions. The work is to obtain
this congruence with a dimension threshold independent of the number
of new primes. Section~\ref{sec:binary-algebra} does so after the new
local blocks have been eliminated from the Selmer complex: a clearing
factor for the remaining bounded complex makes finitely many coefficient
congruences sufficient. Only that clearing factor needs a uniform pole
bound; the numerator may involve the whole prime set.

\subsection{Transfer of a central value}

For a twist with nonzero central value, put
\begin{equation}\label{eq:transfer-normalized-value}
                  D(h)=v(\mathcal L(h))-2w(h)-s(h).
\end{equation}
The subtraction records the local contributions of the ramified primes
and the finite Selmer contribution. In the cyclotomic proof, Kato's class
and the invariant line at the real place form the determinant coordinate
$u_x$. Its central valuation differs from $D(h_x)$ by a fixed bounded
amount. Thus an upper bound on $D$ is exactly the arithmetic input to
the preceding contradiction.

\begin{theorem}[Cyclotomic missing vertex]
\label{thm:cyclotomic-interpolation}
There is a constant $C_E$, depending only on $E$, with the following
property. Fix an odd fundamental discriminant $h_0$, prime to $2N$, with
$c_2(E^{(h_0)})=0$. Suppose that for arbitrarily
large $b$ there is a family
\begin{equation}\label{eq:cy:cube}
 h_x=h_0\prod_{q\in\mathcal Q}(q^*)^{\lambda_q(x)},
 \qquad x\in\F_2^b,\qquad \lambda_q\in(\F_2^b)^*,
\end{equation}
where $\mathcal Q$ consists of distinct primes outside $2Nh_0$,
$q^*=(-1/q)q$, and every $h_x$ has the same sign and the same local
squareclasses at $2N$. Assume that a constant $B$, independent of $b$,
satisfies
\begin{equation}\label{eq:cy:punctured-bound}
 L(E^{(h_x)},1)\ne0,\qquad D(h_x)\le B\qquad(x\ne0).
\end{equation}
If $E[2]$ is irreducible, assume also that some $L$, independent of $b$,
satisfies
\begin{equation}\label{eq:cy:prime-bound}
 v_2(q^2-1)\le L\qquad(q\in\mathcal Q).
\end{equation}
Then $L(E^{(h_0)},1)\ne0$ and
\begin{equation}\label{eq:cy:transferred-bound}
 D(h_0)\le B+C_E.
\end{equation}
The necessary dimension may depend on $h_0$, its finite Selmer group,
and the bound in Equation~\eqref{eq:cy:prime-bound}; $C_E$ does not.
\end{theorem}

The final constant $C_E$ bounds the specialization error. The finite
Selmer group of the base and, in the irreducible case, the bound on
$v_2(q^2-1)$ affect how large a cube is needed, but not this transferred
error. We shall construct the nonzero vertices so that
$v(\mathcal L(h_x))\le2w(h_x)+C$; since $s(h_x)\ge0$, this supplies
the hypothesis on $D(h_x)$.

\subsection{Transfer of a Heegner point}
\label{sec:h:statements}

In corank one the coordinate comes from a Heegner point. The family
must also carry a companion twist whose central value is nonzero.
We first specify the two Heegner constructions whose indices enter
the estimates.

Let $K=\Q(\sqrt{k})$, where $k<0$ is an odd fundamental discriminant and
every prime dividing $2N$ splits in $K$. We omit the two fields with extra
units and, if $E$ has complex multiplication, its CM field. Fix a nonzero modular parametrization $X_0(N)\to E$, translated so
that the cusp $\infty$ maps to the origin, and compatible orientations
of the Heegner points $y_c\in E(H_c)$, where $H_c$ is the ring class
field of conductor $c$. The Manin--Drinfeld theorem makes the images of
cusp differences torsion \citep[Section~1]{Drinfeld1973}. Since the
modular curve and parametrization are fixed, one integer kills all
these images and removes cuspidal ambiguities.

If $\chi_e|_{G_K}$ is a primitive ring class character of conductor $c$,
write $P_c(e)$ for the character \emph{sum} of $y_c$. Equivalently, twist
over $H_c$ and trace to $E^{(e)}(K)$. If this latter group has rank one,
$j(P_c(e))$ is the valuation of its index on the free rank-one lattice.
Passing between this lattice and the rank-one lattice of the contributing
curve among $E^{(e)}(\Q)$ and $E^{(ek)}(\Q)$ changes the index by a bounded
amount: restriction and addition or subtraction of conjugates have
composites multiplication by $2$.

We use two choices of data:
\begin{align*}
 \text{genus data:}\quad& k=h h_*,\quad (h,h_*)=1,
                         \quad e=h_*,\quad c=1,
                         \quad h_*\text{ fixed};\\
 \text{ring data:}\quad& (k,h)=1,\quad e=h,\quad c=|h|.
\end{align*}
The character in genus data is unramified. In ring data, restriction at
each $q\mid h$ has conductor exponent one: this follows by pulling the
quadratic character back under the local norm. Thus these are the asserted
primitive ring class conductors.

\begin{proposition}[Comparison of the two Heegner constructions]
\label{prop:heegner-comparison}
For either set of data,
\begin{equation}\label{eq:h:height}
 \widehat h(P_c(e))
 =C_E\,c\sqrt{|k|}\,
       \left[L(E^{(e)},s)L(E^{(ek)},s)\right]'_{s=1},
\end{equation}
with a fixed nonzero constant determined by the parametrization and height
conventions. In particular, a simple zero of the product gives total
Mordell--Weil rank one and finite Shafarevich--Tate groups.

Suppose $h$ has analytic rank one, and use genus data with a fixed partner
$h_*$ of analytic rank zero and ring data with a possibly varying $k$,
where $L(E^{(hk)},1)\ne0$. Then
\begin{equation}\label{eq:h:index-comparison}
 2j(P_{|h|}(h))-2j(P_1(h_*))
       =v(\mathcal L(hk))+O_{E,h_*}(1).
\end{equation}
Here the two points are formed in their respective quadratic fields.
If instead $h_*$ has analytic rank one and $L(E^{(h)},1)\ne0$, genus
data give
\begin{equation}\label{eq:h:reverse-comparison}
 2j(P_1(h_*))=v(\mathcal L(h))+O_{E,h_*}(1).
\end{equation}
\end{proposition}


The genus construction varies the quadratic field $\Q(\sqrt{hh_*})$
while holding $h_*$ fixed. It is the construction detected by the odd
coefficients below. Interpolation instead uses the ring construction:
within one binary family the field $K$ is common to every vertex.
Equation~\eqref{eq:h:index-comparison} converts a genus-index estimate
and an even central-value estimate into the ring-index estimate needed
for transfer. The field $K$ may change when a new family is chosen.

\begin{theorem}[Heegner interpolation]\label{thm:heegner-interpolation}
Fix $E$ and an initial discriminant $h_0$. Consider families
$h_x$, $x\in\F_2^b$, obtained from $h_0$ by adjoining new prime factors
with linear activation functions on $\F_2^b$. No old factor is removed;
the signed factor belonging to an odd prime $q$ is
$q^*=(-1/q)q$. Suppose the local twisting classes at $2N$ are constant.
For each family let $k$ define an imaginary field satisfying the standing
Heegner conditions, with $(k,h_x)=1$ for every $x$. Each new prime must
either split in $K$ or satisfy $t_q=0$.

Assume that the following bounds are fixed independently of $b$:
\begin{enumerate}[label=\textup{(\roman*)}]
 \item a bound for $\omega(k)$;
 \item a bound for the length of the finite quotient of
       $\Sel_{2^\infty}(E^{(h_0)}/K)$ by its divisible subgroup;
 \item a constant $B$.
\end{enumerate}
Assume this base Selmer group has corank one. Suppose, for every $x\ne0$,
the product $L(E^{(h_x)},s)L(E^{(h_xk)},s)$ has a simple zero and
\begin{equation}\label{eq:h:interpolation-hypothesis}
 2j(P_{|h_x|}(h_x))-4w(h_x)-s(h_x)-s(h_xk)\le B.
\end{equation}
Then, for sufficiently large $b$, the base product has a simple zero.
The dimension threshold depends only on the fixed base and the displayed
bounds, not on the number or sizes of the new primes. The field $K$ may
vary with the family.
\end{theorem}

In these ring data, the relevant normalized quantity is
\[
       2j(P_{|h|}(h))-4w(h)-s(h)-s(hk).
\]
We use the theorem to prove nonvanishing at the base. In our applications,
the genus-index and even central-value estimates will give
$2j(P_{|h_x|}(h_x))\le4w(h_x)+C$ at every nonzero vertex. Subtracting
the nonnegative Sha lengths supplies the required upper bound.

The finite-quotient condition allows two uses of the theorem. The
rational-two-torsion argument fixes $K$ before increasing $b$. The
irreducible argument lets $K$ vary and verifies the following criterion.

\begin{lemma}\label{lem:h:base-selmer-bound}
The finite-quotient bound in Theorem~\ref{thm:heegner-interpolation} holds
if $K$ belongs to a fixed finite collection. It also holds if
$E^{(h_0)}$ is fixed with Selmer corank one, the curves
$E^{(h_0k)}$ have analytic rank zero with bounded $s(h_0k)$, and
$\omega(k)$ is bounded.
\end{lemma}

\subsection{Lower bounds and the coefficient constructions}

The transfer criteria ask for upper bounds at a well-chosen set of
twists. To find those twists, we first need lower bounds on all the
normalized coefficients under consideration. In even sign a least
normalized valuation can then be selected, and a coefficient attaining
it gives both analytic nonvanishing and an upper bound with a constant
independent of the prime support. In odd sign the coefficients come at
increasing finite precision: the minimum is defined through bounded-weight
sequences, and each use is realized on a finite collection of indices.

For the even coefficient construction,
Section~\ref{cs:coefficient-systems} relates twice a coefficient valuation
to $v(\mathcal L(h))$ up to a fixed constant. For the odd construction it
relates the coefficient valuation to the genus index. The following
three lower estimates supply the normalizations, with different roles
in the two residual cases.

When $E[2]$ is reducible, the cyclotomic argument supplies the even
lower estimate directly. By Lemma~\ref{pr:residual-cases}, the
representation is then an extension of two trivial modules.

\begin{theorem}[Cyclotomic lower bound]\label{thm:cyclotomic-bound}
Suppose that $E[2]$ is an extension of two trivial $\F_2$-modules.
There is a constant $C_E$ such that every odd fundamental discriminant
$h$ prime to $2N$ with
$L(E^{(h)},1)\ne0$ satisfies
\begin{equation}\label{eq:cy:lower-bound}
 v(\mathcal L(h))\ge 2w(h)+s(h)-C_E.
\end{equation}
The constant is independent of the number and sizes of the primes
that divide $h$.
\end{theorem}

When $E[2]$ is irreducible, the genus construction supplies the lower
estimate for either sign. With an analytic-rank-zero fixed partner,
it bounds the odd genus index. With an analytic-rank-one fixed
partner, Equation~\eqref{eq:h:reverse-comparison} turns it into the
even central-value bound.

\begin{theorem}[Genus bound]\label{thm:heegner-genus-bound}
In genus data, suppose the product in
Equation~\eqref{eq:h:height} has a simple zero and $E[2]$ is irreducible.
Then, uniformly in $h$,
\begin{equation}\label{eq:h:genus-bound}
 2j(P_1(h_*))\ge 2w(h)+s(h)-O_{E,h_*}(1).
\end{equation}
In proving the uniform estimate, one may omit the finitely many $h$ for
which $K$ meets $\Q(E[2])$ nontrivially.
\end{theorem}

The remaining case is the odd coefficient construction with rational
two-torsion. Here the genus estimate just stated is unavailable. We use
instead the following ring estimate, whose residual representation is
unrestricted.

\begin{theorem}[Ring bound]\label{thm:heegner-ring-bound}
In ring data, suppose the product in
Equation~\eqref{eq:h:height} has a simple zero, every $q\mid h$ splits in
$K$, and $\omega(k)\le a$, where $\omega$ counts distinct prime factors.
Then
\begin{equation}\label{eq:h:ring-bound}
 2j(P_{|h|}(h))\ge
       4w(h)+s(h)+s(hk)-O_{E,a}(1).
\end{equation}
There is no irreducibility assumption on $E[2]$.
\end{theorem}

The passage from this ring estimate to an odd genus estimate is part
of Proposition~\ref{prop:rational-odd-bound}. The even coefficient construction
first provides, for each already rank-one twist $h$ in the positive
filter, a companion parameter $k'$ satisfying the ring-bound hypotheses:
it is negative, odd fundamental, coprime to $2Nh$, a local square at
every prime of $2Nh$, and defines a nonexceptional imaginary field.
Its prime count is bounded, and the companion obeys
$L(E^{(hk')},1)\ne0$ and
\[
       v(\mathcal L(hk'))\le2w(hk')+s(hk')+C.
\]
Substitute this upper bound into
Equation~\eqref{eq:h:index-comparison} and the ring lower bound.
Since $w(hk')=w(h)+w(k')$ and $w(k')$ is bounded, the result is
\[
                   2j(P_1(h_*))\ge2w(h)+s(h)-C'.
\]
Only then is the odd coefficient system normalized in the reducible
branch. Thus the lower-bound dependencies are
\[
 \begin{array}{c@{\quad\longrightarrow\quad}l}
 \text{irreducible genus bound}&\text{even and odd coefficient bounds},\\[2pt]
 \text{reducible cyclotomic bound}&\text{even coefficients, then bounded companions},\\[2pt]
 \text{ring bound and those companions}&\text{reducible odd coefficient bound}.
 \end{array}
\]
The reducible odd normalization is therefore available after the even
companion construction. The all-split hypothesis belongs to the ring
lower bound used here. The Heegner transfer criterion retains its stated
alternative for new primes with $t_q=0$.

\subsection{Auxiliary twists in prescribed local classes}

The coefficient constructions also require an auxiliary twist with
specified local classes. For even sign the number of its prime factors
must remain bounded as the local conditions vary. The following
nonvanishing result supplies both kinds of partner; its proof is given
in Section~\ref{sec:auxiliary-nonvanishing}.

\begin{proposition}[Auxiliary nonvanishing]
\label{prop:rough-nonvanishing}\label{nv:rough}\label{nv:local-witness}
Let $E/\Q$ have conductor $N$, and let $S$ be a finite set of primes
containing the prime divisors of $2N$. Prescribe a real sign
$\sigma\in\{1,-1\}$ and a unit squareclass
$c_p\in\Z_p^\times/\Z_p^{\times2}$ for each $p\in S$.
Suppose that these data are realized by an odd fundamental discriminant,
and that twists in this filter have root number $\varepsilon$.
For every finite set $B$ of primes there are odd fundamental
discriminants $D$ of arbitrarily large absolute value such that
\[
 \operatorname{sign}(D)=\sigma,\qquad [D]_p=c_p\quad(p\in S),\qquad
 p\nmid D\quad(p\in B),
\]
Moreover:
\begin{enumerate}[label=\textup{(\roman*)}]
\item If $\varepsilon=1$, they can be chosen with
$L(E^{(D)},1)\ne0$ and $\omega(|D|)\le5$.
\item If $\varepsilon=-1$, they can be chosen with
$L'(E^{(D)},1)\ne0$, so the central zero is simple.
\end{enumerate}
The bound five in \textup{(i)} is independent of $E,S,B$ and the
prescribed classes. In either case, every prime factor of $D$ can be
required to exceed any previously fixed bound.
\end{proposition}
